\begin{table}[htbp]
\centering
\caption{Metrik evaluasi sistem penelusuran dan pembangkitan.}
\label{tab:metrik-evaluasi}
\small
\begin{tabular}{@{}p{3.0cm}p{5.5cm}p{1.8cm}p{2.3cm}@{}}
\toprule
\textbf{Metrik} & \textbf{Aspek yang diukur} & \textbf{Rentang} & \textbf{Rujukan} \\
\midrule
\multicolumn{4}{@{}l}{\emph{Mutu penelusuran}} \\[2pt]
Hit@k & Apakah potongan relevan terjangkau di antara $k$ hasil teratas & 0 sampai 1 & \textcite{manning2008} \\[2pt]
\emph{Mean reciprocal rank} & Seberapa awal potongan relevan pertama muncul pada daftar hasil & 0 sampai 1 & \textcite{manning2008} \\[4pt]

\multicolumn{4}{@{}l}{\emph{Mutu pembangkitan}} \\[2pt]
\emph{Faithfulness} & Kesetiaan jawaban terhadap konteks yang diperoleh & 0 sampai 1 & \textcite{es2023} \\[2pt]
\emph{Answer relevancy} & Relevansi jawaban terhadap pertanyaan yang diajukan & 0 sampai 1 & \textcite{es2023} \\[2pt]
\emph{Context precision} & Proporsi konteks yang diambil yang benar-benar relevan & 0 sampai 1 & \textcite{es2023} \\[2pt]
\emph{Context recall} & Proporsi konteks relevan yang berhasil diambil & 0 sampai 1 & \textcite{es2023} \\[4pt]

\multicolumn{4}{@{}l}{\emph{Mutu dan keamanan prediksi}} \\[2pt]
RMSE dan MAE & Besar galat prediksi terhadap nilai aktual & mg/dL & --- \\[2pt]
\emph{Clarke Error Grid} & Konsekuensi klinis dari galat, dipetakan ke zona A sampai E & zona A--E & \textcite{clarke1987} \\[2pt]
Sensitivitas dan PPV & Kemampuan mengenali kejadian hipoglikemia dan keandalan peringatannya & 0 sampai 1 & --- \\[2pt]
Cakupan konformal & Proporsi nilai aktual yang tercakup interval prediksi & 0 sampai 1 & \textcite{angelopoulos2021} \\
\bottomrule
\end{tabular}

\vspace{2pt}
\footnotesize Rumusan matematis setiap metrik disajikan pada Bab~\ref{chap:evaluasi} agar dapat disertai penjelasan komponennya.\normalsize
\end{table}
